\documentclass[aps,prl,reprint]{revtex4-2}

\usepackage[utf8]{inputenc}
\usepackage{graphicx}
\usepackage{amsmath}
\usepackage{amssymb}
\usepackage{caption}
\usepackage{subcaption}

\begin{document}
	
\title{Project Quantum Go: Ising Feature Maps and Photonic Graph Features}

\author{M. A. Mercado Sánchez$^{1}$, L. Jiménez$^{2}$, and B. Escalante Ramírez$^{1}$}

\affiliation{
$^{1}$Departamento de Procesamiento de Señales, División de Ingeniería Eléctrica, Facultad de Ingeniería, Universidad Nacional Autónoma de México (UNAM) \\
$^{2}$Facultad de Ciencias, Universidad Nacional Autónoma de México (UNAM)
}

\begin{abstract}
We study the game of Go as a source of features for machine learning tasks (e.g., winner prediction, strategic move evaluation) using two complementary approaches: (i) a symmetric classical Ising model translated to QUBO for compatibility with adiabatic quantum annealers, and (ii) photonic-inspired graph features via Gaussian Boson Sampling (GBS) over a chain-configuration graph (CFG) of the board. We clarify hardware constraints inherent to quantum annealers (which solve classical, not quantum, Ising/QUBO problems), present a symmetric reformulation using auxiliary bits to remove color bias present in baseline models, and outline how GBS-derived features capture higher-order group structure on the Go board. We validate our approach on 6,000 professional game records. Pending experimental work (execution on real D-Wave hardware and complete GBS feature validation) is explicitly marked as planned work.
\end{abstract}

\maketitle

\section{Introduction}

\subsection{Motivation and scope}

The ancient game of Go presents a grand challenge for computational analysis due to its vast state space ($\sim 10^{170}$ legal positions) and subtle strategic depth. Traditional approaches combine Monte Carlo Tree Search (MCTS) with deep neural networks trained on expert games. Here, we explore an alternative paradigm: extracting \textit{physics-inspired features} from Go positions using quantum computing techniques.

Our goal is to derive interpretable energy-based and graph-theoretic features that can feed downstream learning models for tasks such as predicting game outcomes or identifying critical moves. We pursue two complementary lines of investigation:

\begin{enumerate}
	\item \textbf{Ising/QUBO formulation:} A hardware-aligned classical Ising model that can be mapped to Quadratic Unconstrained Binary Optimization (QUBO) and, in principle, solved on quantum annealers like D-Wave systems.
	
	\item \textbf{Photonic graph features:} A Gaussian Boson Sampling (GBS)-inspired approach that maps the board's chain-configuration graph to an interferometer, extracting subgraph statistics and correlation patterns.
\end{enumerate}

\subsection{Clarification on quantum vs. classical}

A critical distinction must be made at the outset: quantum annealers (such as D-Wave systems) \textit{do not} execute arbitrary quantum Hamiltonians with non-diagonal terms (e.g., $X$ and $Z$ cross-terms). Instead, they find the ground state of \textit{classical} Ising Hamiltonians, which are diagonal in the computational ($Z$) basis:

\begin{equation}
H_{\text{Ising}} = \sum_i h_i \sigma^z_i + \sum_{i<j} J_{ij} \sigma^z_i \sigma^z_j
\end{equation}

This is equivalent to solving a QUBO problem over binary variables $x_i \in \{0,1\}$. The quantum annealing process uses quantum fluctuations (tunneling) to explore the energy landscape more efficiently than classical thermal annealing, but the target Hamiltonian itself is classical.

In this work:
\begin{itemize}
	\item The \textbf{quantum Ising Hamiltonian} with $X$-$Z$ terms (Section~\ref{sec:quantum_ising}) was explored for \textit{pedagogical and simulation purposes only} using PennyLane. It provides intuition about time evolution and amplitude flows but \textit{cannot be directly executed on D-Wave hardware}.
	
	\item The \textbf{symmetric classical Ising model} (Section~\ref{sec:symmetric_ising}) is the hardware-compatible formulation, designed to be converted to QUBO and embedded on annealer topologies.
\end{itemize}

We did not execute experiments on real D-Wave hardware due to access limitations during the project period; all annealing results presented are based on simulations.

\section{Classical Ising Models for Go}

\subsection{Baseline: The Atomic-Go model}

The Atomic-Go model, proposed by Alvarado~\cite{alvarado2020}, provides a straightforward adaptation of the 2D Ising model to the game of Go. Each position on the board is associated with a spin variable:

\begin{equation}
x_i \in \{-1, +1, 0\}
\end{equation}

where $x_i = +1$ represents a white stone, $x_i = -1$ represents a black stone, and $x_i = 0$ represents an empty intersection. The local energy for a stone at position $i$ is computed as:

\begin{equation}
s_i = x_i \sum_{j \in \mathcal{N}(i)} x_j
\end{equation}

where $\mathcal{N}(i)$ denotes the set of orthogonally adjacent positions (neighbors). The total Hamiltonian is:

\begin{equation}
H_{\text{Atomic}}(t) = -\sum_i s_i^{(t)} - \mu \sum_i h_i^{(t)} x_i^{(t)}
\label{eq:atomic}
\end{equation}

where $h_i^{(t)}$ represents the number of liberties (empty neighbors) of the stone at position $i$, and $\mu \geq 0$ is a parameter controlling the importance of liberties.

\subsubsection{Limitations of the Atomic model}

While computationally simple, the Atomic-Go model suffers from several critical limitations:

\begin{enumerate}
	\item \textbf{Asymmetry:} The energy is color-dependent in a biased way. A black stone ($x_i = -1$) surrounded by white stones ($x_j = +1$) has different energy than a white stone surrounded by black stones, even though the strategic situation is symmetric.
	
	\item \textbf{Lack of constraints:} The model does not encode Go rules such as captures, ko (repetition), or suicide (placing a stone with zero liberties).
	
	\item \textbf{Quadratic form:} While the model is quadratic (suitable for QUBO), it treats all neighbor interactions identically, ignoring tactical structures like eyes, connections, and cutting points.
\end{enumerate}

Figure~\ref{fig:go_and_energy} shows a professional game position and its corresponding energy map computed using the Atomic model. Blue regions (negative energy) correspond to black-dominated areas, red regions (positive energy) to white-dominated areas, and neutral (near-zero energy) regions represent contested or empty territories.

\begin{figure}[h]
	\centering
	\begin{subfigure}[b]{0.45\textwidth}
		\centering
		% \includegraphics[width=\textwidth]{../data/assets/articulo/Tablero.jpg}
		\textit{[Game position image: Park Jungwhan vs. Zhao Chenyu]}
		\caption{Game position: Park Jungwhan vs. Zhao Chenyu}
		\label{fig:board}
	\end{subfigure}
	\hfill
	\begin{subfigure}[b]{0.45\textwidth}
		\centering
		% \includegraphics[width=\textwidth]{../data/assets/articulo/Mapa_energia.jpg}
		\textit{[Atomic-Go energy map image]}
		\caption{Atomic-Go energy map}
		\label{fig:energy_map}
	\end{subfigure}
	\caption{Comparison of a Go board state and its corresponding Atomic-Go energy map. Energies range from $-4$ (blue, black-dominated) to $+4$ (red, white-dominated), with $0$ (neutral) corresponding to white in the color scale.}
	\label{fig:go_and_energy}
\end{figure}

\subsection{Symmetric Ising model with auxiliary bits}
\label{sec:symmetric_ising}

To address the asymmetry problem and prepare for QUBO conversion, we reformulate the model using \textit{auxiliary binary variables}. The key idea is to separate ``color'' from ``occupation'' using two bits per position:

\begin{itemize}
	\item $p_i \in \{0,1\}$: indicates presence of a \textit{positive} (white) stone
	\item $n_i \in \{0,1\}$: indicates presence of a \textit{negative} (black) stone
\end{itemize}

The constraint $p_i + n_i \leq 1$ ensures that at most one stone color occupies position $i$ (enforced via energy penalties). The spin variable is reconstructed as:

\begin{equation}
x_i = p_i - n_i \in \{-1, 0, +1\}
\end{equation}

\subsubsection{Target symmetric energy table}

For two adjacent positions $i$ and $j$, we define a symmetric 2-body energy table that treats black and white identically. Table~\ref{tab:symmetric_energy} shows the desired energy contributions:

\begin{table}[h]
	\centering
	\caption{Symmetric 2-body energy table for adjacent positions. $B$ = black stone, $W$ = white stone, $E$ = empty. Energy is computed from the perspective of color-matching and liberties.}
	\label{tab:symmetric_energy}
	\begin{tabular}{|c|c|c|}
		\hline
		Position $i$ & Position $j$ & Energy contribution \\
		\hline
		$B$ & $B$ & $-J$ (friendly neighbor) \\
		$W$ & $W$ & $-J$ (friendly neighbor) \\
		$B$ & $W$ & $+J$ (opposing neighbor) \\
		$W$ & $B$ & $+J$ (opposing neighbor) \\
		$B$ & $E$ & $-\mu$ (liberty bonus) \\
		$W$ & $E$ & $-\mu$ (liberty bonus) \\
		$E$ & $B$ & $-\mu$ (liberty bonus) \\
		$E$ & $W$ & $-\mu$ (liberty bonus) \\
		$E$ & $E$ & $0$ (no interaction) \\
		\hline
	\end{tabular}
\end{table}

This table is color-symmetric by construction: swapping all black and white stones produces the same energy landscape (up to a global sign, which does not affect optimization).

\subsubsection{QUBO conversion and penalties}

To enforce the constraint $p_i + n_i \leq 1$ and implement the energy table, we add penalty terms to the Hamiltonian:

\begin{equation}
H_{\text{penalty}} = \lambda \sum_i p_i n_i
\end{equation}

where $\lambda \gg J, \mu$ is a large penalty coefficient that makes configurations with $p_i = n_i = 1$ energetically forbidden.

The full Hamiltonian in $(p, n)$ variables becomes:

\begin{align}
H_{\text{symmetric}} &= \sum_{\langle i,j \rangle} \left[ J_{ij}(p_i, n_i, p_j, n_j) + \mu L_{ij}(p_i, n_i, p_j, n_j) \right] \nonumber \\
&\quad + \lambda \sum_i p_i n_i + H_{\text{constraints}}
\label{eq:symmetric_hamiltonian}
\end{align}

where:
\begin{itemize}
	\item $J_{ij}(\cdot)$ implements the interaction terms from Table~\ref{tab:symmetric_energy}
	\item $L_{ij}(\cdot)$ implements liberty bonuses
	\item $H_{\text{constraints}}$ encodes additional Go rules (captures, ko, suicide) as soft constraints (penalties for rule violations)
\end{itemize}

The explicit polynomial expansion of $J_{ij}$ and $L_{ij}$ in terms of $p_i, n_i, p_j, n_j$ is quadratic by design. For example, the ``friendly neighbor'' term for black stones can be written as:

\begin{equation}
E_{\text{BB}} = -J \cdot n_i n_j
\end{equation}

and similarly for white-white interactions. The opposing-color term involves products like $p_i n_j$ and $n_i p_j$.

\subsubsection{Handling fixed board positions}

A critical practical issue arises when applying this model to mid-game positions: some board locations already contain fixed stones (from previous moves), while others are empty and represent the search space.

Quantum annealers do not natively support ``pinning'' qubits to fixed values. We address this via:

\begin{enumerate}
	\item \textbf{Strong biases:} For a position $i$ known to contain a black stone, we add large negative biases to $n_i$:
	\begin{equation}
	h_{n_i} = -\Lambda, \quad h_{p_i} = +\Lambda
	\end{equation}
	where $\Lambda \gg J, \mu, \lambda$ makes it energetically almost impossible for the annealer to flip these bits.
	
	\item \textbf{Problem reduction:} Alternatively, we can algebraically eliminate fixed variables, solving a reduced QUBO only over empty positions. This requires recomputing neighbor interactions and is more complex but uses fewer qubits.
\end{enumerate}

In our current implementation, we use the strong-bias approach with $\Lambda \approx 100 \times J$.

\subsection{Embedding on annealer hardware}

D-Wave annealers have specific connectivity topologies (Pegasus, Zephyr). Our QUBO problem must be \textit{embedded} onto the hardware graph, potentially requiring:

\begin{itemize}
	\item \textbf{Chain construction:} Logical qubits are represented by chains of physical qubits with strong ferromagnetic coupling.
	\item \textbf{Coefficient scaling:} The QUBO coefficients $(h_i, J_{ij})$ must fit within hardware precision ranges (typically $\sim [-2, +2]$ in normalized units).
\end{itemize}

\textbf{Zephyr architecture} (available in Advantage2 systems) has higher connectivity ($\sim 20$ edges per qubit) and \textit{programmable couplers} that can natively represent certain cubic terms, reducing the need for auxiliary variables in some cases. However, our current symmetric model already uses auxiliaries for conceptual clarity and color symmetry.

\textbf{Status:} Embedding code is implemented using the D-Wave Ocean SDK's \texttt{minorminer} tool. Coefficient scaling and chain strength tuning are pending calibration on real hardware.

\section{Quantum Ising Model (Simulation Only)}
\label{sec:quantum_ising}

For pedagogical completeness and to explore quantum time evolution, we also investigated a \textit{quantum} Ising Hamiltonian that includes non-diagonal terms:

\begin{equation}
H_{\text{MIG}} = -J \sum_{\langle i,j \rangle} \sigma^z_i \sigma^z_j - h \sum_{\langle i,j \rangle} \sigma^x_i \otimes (H \sigma^x_j)
\label{eq:quantum_hamiltonian}
\end{equation}

where:
\begin{itemize}
	\item $\sigma^z_i$ and $\sigma^x_i$ are Pauli Z and X operators on qubit $i$
	\item $H$ is the Hadamard gate: $H = \frac{1}{\sqrt{2}}(X + Z)$
	\item $\langle i,j \rangle$ denotes pairs of adjacent board positions
\end{itemize}

\subsection{Interpretation}

\begin{itemize}
	\item \textbf{First term} ($\sigma^z_i \sigma^z_j$): Classical Ising interaction between neighboring qubits. Favors alignment (ferromagnetic if $J > 0$, antiferromagnetic if $J < 0$).
	
	\item \textbf{Second term} ($\sigma^x_i \otimes H\sigma^x_j$): Transverse field with Hadamard modulation. Induces quantum fluctuations and entanglement. The Hadamard gate on qubit $j$ creates superposition, representing the ``uncertainty'' of whether position $j$ will be occupied by black ($|0\rangle$) or white ($|1\rangle$).
\end{itemize}

\subsection{Qubit encoding}

Each board position is mapped to a qubit with basis states:
\begin{itemize}
	\item $|0\rangle$: black stone
	\item $|1\rangle$: white stone
	\item Superposition $\alpha|0\rangle + \beta|1\rangle$: empty position (to be determined)
\end{itemize}

\subsection{Time evolution simulation}

We simulate the time evolution of this Hamiltonian on small boards (e.g., $5 \times 5$, 25 qubits) using PennyLane. The initial state is prepared as a uniform superposition:

\begin{equation}
|\psi(0)\rangle = \bigotimes_{i=1}^{N} \frac{1}{\sqrt{2}}(|0\rangle_i + |1\rangle_i)
\end{equation}

The system evolves under the Schrödinger equation:

\begin{equation}
|\psi(t)\rangle = e^{-iH_{\text{MIG}}t/\hbar} |\psi(0)\rangle
\end{equation}

We compute the expectation value of the energy:

\begin{equation}
\langle H \rangle(t) = \langle \psi(t) | H_{\text{MIG}} | \psi(t) \rangle
\end{equation}

and visualize how the amplitude distribution evolves over time. This provides insight into which board configurations are energetically favored under quantum dynamics.

\subsection{Limitations and clarification}

\textbf{This Hamiltonian cannot be executed on D-Wave annealers.} D-Wave systems do not support arbitrary time-dependent Hamiltonians or gate-based quantum circuits. The MIG Hamiltonian is purely a theoretical/simulation construct used to:

\begin{enumerate}
	\item Understand how quantum superposition and entanglement might influence board evaluation.
	\item Provide a conceptual bridge to quantum feature maps in quantum machine learning.
	\item Motivate the search for quantum-inspired classical algorithms (which is what the symmetric Ising + QUBO approach represents).
\end{enumerate}

All results in Section~\ref{sec:quantum_ising} are based on exact state vector simulation using PennyLane on classical hardware (CPUs/GPUs). No quantum hardware was used.

\section{Photonic Graph Features via Gaussian Boson Sampling}
\label{sec:gbs}

In addition to energy-based features from Ising models, we explore a complementary approach based on \textit{graph-theoretic features} extracted via Gaussian Boson Sampling (GBS), a photonic quantum computing paradigm.

\subsection{Chain-Configuration Graph (CFG)}

We represent a Go position as a graph $G = (V, E)$ where:
\begin{itemize}
	\item \textbf{Nodes} ($V$): Each node represents a connected group of same-color stones (a ``chain'' in Go terminology). Isolated stones are single-node groups.
	\item \textbf{Edges} ($E$): Edges connect nodes (chains) that:
	\begin{itemize}
		\item Share liberties (adjacent empty points)
		\item Are spatially close (e.g., within Manhattan distance 2)
		\item Belong to opposing colors and are in tactical contact
	\end{itemize}
	\item \textbf{Edge weights}: Encode strategic importance (e.g., number of shared liberties, cutting points, potential captures).
\end{itemize}

The adjacency matrix $\mathbf{A}$ of this graph is computed from the board state using the parser in \texttt{src/go\_game\_engine.py}.

\subsection{GBS feature extraction}

GBS is a quantum sampling task where photons are injected into a linear optical interferometer and detected at output modes. The interferometer is programmed by a symmetric matrix (related to $\mathbf{A}$ via the Takagi decomposition). The output samples correspond to subsets of nodes (subgraphs) with probabilities influenced by the graph's structure.

For a Go position, we:
\begin{enumerate}
	\item Compute the CFG adjacency matrix $\mathbf{A}$.
	\item Perform the Takagi decomposition: $\mathbf{A} = \mathbf{U} \mathbf{\Lambda} \mathbf{U}^T$, where $\mathbf{\Lambda}$ is diagonal (singular values) and $\mathbf{U}$ is unitary.
	\item Program a (simulated or real) GBS device with parameters derived from $\mathbf{U}$ and $\mathbf{\Lambda}$.
	\item Collect samples: each sample is a subset $S \subseteq V$ representing a subgraph.
	\item Extract statistical features from the sample distribution.
\end{enumerate}

\subsection{Features extracted}

From $N$ samples $\{S_1, S_2, \ldots, S_N\}$, we compute:

\begin{itemize}
	\item \textbf{Mean subgraph size}: $\langle |S| \rangle = \frac{1}{N}\sum_{k=1}^N |S_k|$
	\item \textbf{Maximum clique size}: Largest fully connected subgraph sampled.
	\item \textbf{Node concentration}: Frequency of each node appearing in samples (which chains are ``important'').
	\item \textbf{Sample entropy}: $H = -\sum_{S} p(S) \log p(S)$, where $p(S)$ is the empirical probability of sampling subgraph $S$. High entropy indicates positional ambiguity; low entropy indicates a decisive configuration.
	\item \textbf{Correlation matrix}: $C_{ij} = P(\text{nodes } i \text{ and } j \text{ both in sample})$. Reveals which chains are tactically linked.
	\item \textbf{Temporal similarity}: For successive moves $t$ and $t+1$, compute the graph edit distance or Jaccard similarity between $G^{(t)}$ and $G^{(t+1)}$.
\end{itemize}

\subsection{Motivation and advantages}

Unlike the Ising energy (which is a scalar), GBS features are \textit{multi-dimensional} and capture:
\begin{itemize}
	\item Community structure (which groups of stones form tactical units)
	\item Robustness of connections (high-frequency subgraphs are stable structures)
	\item Ambiguity and tension (high entropy suggests many plausible continuations)
\end{itemize}

These features are complementary to energy maps and may be particularly useful for:
\begin{itemize}
	\item Predicting game outcome (decisive vs. close games)
	\item Identifying critical moments (sudden entropy drops indicate forcing sequences)
	\item Evaluating positional complexity
\end{itemize}

\subsection{Current implementation status}

\begin{itemize}
	\item \textbf{CFG construction}: Fully implemented in Python. Tested on 6,000 professional SGF game records.
	\item \textbf{GBS simulation}: Implemented using the Strawberry Fields library (Xanadu). Samples are generated via the Gaussian backend (exact Hafnian computation for small graphs).
	\item \textbf{Feature extraction}: Code exists for computing all features listed above. Example outputs (entropy curves, correlation heatmaps) are included in the repository notebooks.
	\item \textbf{Real hardware execution}: Not yet performed. Xanadu's Borealis photonic processor is publicly accessible but was not used during this project due to time constraints. Future work will validate simulation results against real photonic hardware.
\end{itemize}

\section{Data Processing Pipeline}

\subsection{SGF game corpus}

We obtained 6,000 professional Go game records in Smart Game Format (SGF) from the Gofiku.com database. These games span:
\begin{itemize}
	\item Board sizes: primarily $19 \times 19$ (standard), with some $9 \times 9$ for testing
	\item Player ranks: amateur dan to professional 9-dan
	\item Game outcomes: recorded as Black win (B+), White win (W+), or draw
\end{itemize}

\subsection{Rules engine}

The parser and rules engine (\texttt{src/go\_game\_engine.py}) implements:
\begin{itemize}
	\item Move validation and execution
	\item Stone capture detection (groups with zero liberties)
	\item Ko detection (simple ko and superko variants)
	\item Suicide rule enforcement (depends on ruleset)
	\item Liberty counting for each group
	\item Board state hashing for repetition detection
\end{itemize}

Each move is processed sequentially, producing a time series of board states.

\subsection{Feature computation}

For each board state $t$:
\begin{enumerate}
	\item \textbf{Atomic-Go energy map}: Computed via Equation~\eqref{eq:atomic}. Stored as an $N \times N$ matrix (where $N$ is board size).
	\item \textbf{Symmetric Ising energy}: (Pending full implementation) Will be computed via Equation~\eqref{eq:symmetric_hamiltonian} once QUBO coefficients are calibrated.
	\item \textbf{CFG adjacency matrix}: Computed from connected components analysis.
	\item \textbf{GBS features}: Extracted from 1,000 samples per position (simulation).
\end{enumerate}

\subsection{Visualization and export}

The visualization module (\texttt{src/go\_visualization.py}) generates:
\begin{itemize}
	\item Static PNG images of board states and energy maps (Figure~\ref{fig:go_and_energy})
	\item Animated GIFs showing time evolution of energy maps throughout a game
	\item Interactive HTML5 game browsers with move-by-move navigation, liberty display, and energy overlays
	\item Correlation heatmaps for GBS node frequencies
\end{itemize}

Documentation for the interactive analysis tool is provided in \texttt{docs/analisis\_interactivo\_partidas\_go.md}.

\section{Experimental Results and Current Status}

\subsection{Atomic-Go energy validation}

We computed Atomic-Go energy maps for all 6,000 games. Preliminary analysis shows:
\begin{itemize}
	\item \textbf{Correlation with outcome}: The final board energy (summed over all positions) has weak correlation with game outcome (Pearson $r \approx 0.3$). This is expected because the Atomic model ignores captures and territory scoring.
	\item \textbf{Energy dynamics}: Energy fluctuates significantly during opening and middle game, then stabilizes in the endgame. Sharp energy changes often correspond to captures or ko fights.
	\item \textbf{Visualization quality}: Energy maps (Figure~\ref{fig:go_and_energy}) are visually interpretable and align with human intuition about ``spheres of influence.''
\end{itemize}

\subsection{Symmetric Ising QUBO}

\begin{itemize}
	\item \textbf{2-qubit toy model}: Validated on a $1 \times 2$ board (two adjacent positions). The QUBO correctly reproduces the target energy table (Table~\ref{tab:symmetric_energy}) for all 9 possible configurations.
	\item \textbf{Full board scaling}: Code exists to generate the full QUBO matrix for any board size. Tested up to $5 \times 5$ (25 positions, 50 binary variables with $p/n$ encoding).
	\item \textbf{Coefficient scaling and embedding}: Pending. We need to tune chain strengths and test embedding success rates on Pegasus/Zephyr topologies.
	\item \textbf{Annealing runs}: Not yet performed (no hardware access). Simulated annealing on classical hardware (simulated annealing via \texttt{neal} sampler) shows qualitatively reasonable solutions but does not demonstrate quantum advantage.
\end{itemize}

\subsection{GBS features}

\begin{itemize}
	\item \textbf{CFG extraction}: Successfully tested on 100 sample games. Average graph size: 15--30 nodes for mid-game $19 \times 19$ positions.
	\item \textbf{Sample entropy}: Computed for selected games. Opening positions show high entropy ($H \approx 3$--$4$ nats), while endgame positions show low entropy ($H \approx 1$--$2$ nats), consistent with the intuition that endgames are more deterministic.
	\item \textbf{Clique sizes}: Maximum cliques of size 5--8 detected, corresponding to solid groups of stones.
	\item \textbf{Aggregated metrics vs. outcome}: Pending. We plan to train a logistic regression model using GBS features to predict game outcome and compare against a baseline using only board statistics (stone count, territory estimate).
\end{itemize}

\subsection{Quantum Hamiltonian simulation}

Time evolution under $H_{\text{MIG}}$ (Equation~\eqref{eq:quantum_hamiltonian}) was simulated for $5 \times 5$ boards. Results show:
\begin{itemize}
	\item Amplitude flows toward low-energy configurations over time.
	\item Entanglement between adjacent qubits grows, then partially decays (due to dephasing in the simulation).
	\item This is purely illustrative; no quantum hardware was involved.
\end{itemize}

\section{Discussion}

\subsection{Advantages of the symmetric Ising formulation}

Compared to the Atomic-Go model, the symmetric formulation:
\begin{enumerate}
	\item \textbf{Removes color bias}: Black and white stones are treated identically.
	\item \textbf{Hardware-compatible}: Can be mapped to QUBO and embedded on D-Wave topology.
	\item \textbf{Extensible}: Additional constraints (captures, ko) can be encoded as penalty terms.
\end{enumerate}

\subsection{Challenges and limitations}

\begin{enumerate}
	\item \textbf{Hyperparameter tuning}: The model has several free parameters ($J, h, \mu, \lambda, \Lambda$) that must be tuned for optimal performance. Grid search over these parameters is computationally expensive.
	
	\item \textbf{Embedding overhead}: Real D-Wave hardware has limited connectivity. Embedding a $19 \times 19$ board may require thousands of physical qubits due to chain construction. Smaller boards ($9 \times 9$) are more feasible.
	
	\item \textbf{Fixed qubits}: The strong-bias approach for fixing stones works but wastes problem capacity. Problem reduction (algebraic elimination) would be more efficient but harder to implement.
	
	\item \textbf{Validation}: Without access to real annealing hardware, we cannot definitively claim quantum advantage or even practical utility. Classical simulated annealing may perform comparably.
\end{enumerate}

\subsection{Complementarity of GBS features}

GBS features:
\begin{itemize}
	\item \textbf{Do not require explicit rule encoding}: The graph structure implicitly captures tactical relationships.
	\item \textbf{Are hardware-specific}: Depend on photonic quantum sampling, not annealers.
	\item \textbf{Provide higher-order information}: Subgraph statistics go beyond pairwise (2-body) interactions.
\end{itemize}

However, GBS features are currently illustrative. Real photonic hardware (Xanadu Borealis) would be needed to validate that quantum sampling provides advantage over classical graph algorithms (e.g., MCMC on the graph).

\section{Future Work}

\subsection{Immediate next steps}

\begin{enumerate}
	\item \textbf{Obtain D-Wave access}: Execute the symmetric Ising QUBO on real Advantage2 hardware (Zephyr topology). Benchmark embedding feasibility and annealing quality.
	
	\item \textbf{Complete GBS validation}: Run Borealis experiments on representative game positions. Compare sample distributions against classical baseline (uniform random sampling, MCMC).
	
	\item \textbf{Train predictive models}: Use energy + GBS features as input to machine learning models (logistic regression, random forest, neural network) to predict game outcomes. Compare against baselines using only board statistics.
	
	\item \textbf{Integrate into interactive UI}: Add an ``energy tab'' to the game browser showing real-time energy maps and GBS features as users navigate moves.
\end{enumerate}

\subsection{Longer-term directions}

\begin{enumerate}
	\item \textbf{Quantum advantage benchmarking}: Rigorously compare quantum annealing / GBS against state-of-the-art classical methods (MCTS, neural network evaluation functions) on standardized datasets (e.g., GoGoD database).
	
	\item \textbf{Incorporate domain knowledge}: Encode additional Go concepts (eyes, seki, sente/gote) into the Hamiltonian or graph structure.
	
	\item \textbf{Quantum-classical hybrid}: Use quantum features as inputs to classical deep learning models (quantum feature kernels).
	
	\item \textbf{Robustness analysis}: Test sensitivity to noisy or incomplete SGF records, and to variations in hyperparameters.
\end{enumerate}

\section{Conclusions}

We have developed a framework for extracting physics-inspired features from Go game positions using two complementary quantum computing paradigms:

\begin{enumerate}
	\item A \textbf{symmetric classical Ising model} that removes color bias, encodes Go rules as constraints, and is compatible with quantum annealers (D-Wave) via QUBO formulation.
	
	\item A \textbf{photonic graph feature approach} based on Gaussian Boson Sampling over chain-configuration graphs, which captures higher-order tactical structure without explicit rule encoding.
\end{enumerate}

We clarified the critical distinction between classical and quantum Hamiltonians in the context of annealing hardware, and honestly reported that no real quantum hardware was used (all results are simulations). The framework is validated on 6,000 professional game records and includes a complete data processing pipeline with visualization tools.

Pending experimental work—execution on D-Wave Advantage2 and Xanadu Borealis, and integration of features into predictive models—will quantify whether these quantum-inspired approaches provide practical advantage over classical baselines. Regardless of outcome, the project demonstrates a rigorous methodology for applying quantum computing concepts to combinatorial game problems.

\begin{acknowledgments}
This work was developed during a postdoctoral fellowship supported by CONAHCYT (CVU 469604) at the Facultad de Ingeniería, Universidad Nacional Autónoma de México (UNAM), under the supervision of Dr. Boris Escalante Ramírez. We thank the Departamento de Procesamiento de Señales for infrastructure support and the UNAM Quantum Computing Working Group for valuable discussions. We acknowledge Xanadu and D-Wave for providing public access to documentation and simulation tools, and the Go community for maintaining open game databases.
\end{acknowledgments}

\begin{thebibliography}{99}
	
	\bibitem{alvarado2020}
	M. Alvarado,
	\textit{Using Deep Belief Networks and the Atomic Go Model to Predict the Winner of Go Games},
	(2020). [Thesis or technical report - update citation]
	
	\bibitem{lucas2014}
	A. Lucas,
	\textit{Ising formulations of many NP problems},
	Frontiers in Physics \textbf{2}, 5 (2014).
	
	\bibitem{dwave_ocean}
	D-Wave Systems,
	\textit{Ocean Software Documentation},
	\texttt{https://docs.ocean.dwavesys.com/}
	
	\bibitem{strawberryfields}
	N. Killoran et al.,
	\textit{Strawberry Fields: A Software Platform for Photonic Quantum Computing},
	Quantum \textbf{3}, 129 (2019).
	
	\bibitem{projectrepo}
	M. A. Mercado Sánchez,
	\textit{Project Quantum Go repository},
	\texttt{https://github.com/ometitlan/Project-Quantum-Go} (2024).
	
\end{thebibliography}

\appendix

\section{Go Rules Formalization}
\label{app:rules}

Go rules (Japanese, Chinese, or Tromp-Taylor variants) can be encoded as penalty terms in the Hamiltonian:

\subsection{Capture constraint}

A stone at position $i$ with zero liberties should be removed (captured). In energy terms, we penalize configurations where:
\begin{equation}
x_i \neq 0 \quad \text{and} \quad h_i = 0
\end{equation}

This can be enforced via:
\begin{equation}
H_{\text{capture}} = \kappa \sum_i (x_i^2)(1 - \mathbb{I}[h_i > 0])
\end{equation}
where $\kappa$ is a penalty weight and $\mathbb{I}[\cdot]$ is an indicator function (in practice, approximated via polynomial constraints).

\subsection{Ko rule}

The ko rule forbids immediate recapture that would return to the previous board state. This requires board state hashing and a penalty that increases energy if the current state matches the state two moves ago.

\subsection{Suicide rule}

Depending on ruleset, placing a stone that immediately has zero liberties may be forbidden. This is handled similarly to the capture constraint.

These constraints significantly increase the QUBO complexity and are subjects of ongoing work.

\end{document}